
Most machine learning datasets are invisible --- not because they lack quality, but because they lack keyword tags that enable search-indexed discovery. This study examines whether keyword tagging, the operationalization of the FAIR F1 (Findability) principle on OpenML, predicts dataset adoption measured by ML task registrations. Analyzing 5,217 datasets with negative binomial (NB-2) regression and decade fixed effects, binary tag presence is associated with 22.6\% more task registrations (IRR=1.2263, 95\% CI [1.1681, 1.2873], p=1.$87\times 10^{-16})$. This effect survives decade fixed effects with only 12.2\% attenuation (attenuation ratio=1.1219) despite extreme era-tag collinearity (Cramér's V=0.823), confirming a structural platform property rather than a temporal artifact. Within tagged datasets (N=2,625), tag count amplifies adoption log-linearly (IRR=1.5332 per log(tag\_count+1), 95\% CI [1.4680, 1.6014], p=1.$28\times 10^{-82})$, with essentially zero decade confounding (attenuation ratio=1.0007). Categorical analysis reveals that OpenML tagging behavior is bimodal --- the 6+ tag tier (IRR=1.2861) is the sole meaningfully distinguishable amplification tier; intermediate bins (1-2 and 3-5 tags) produce statistically indistinguishable effects. These findings constitute the first NB-2 quantification of FAIR F1 keyword tagging on an ML platform, translating aspirational data management principles into actionable, evidence-grounded guidance.

